\subsection{Does a longer output limit increase the value of indexed admission?}\label{sec:cap}
An output limit can affect two parts of a voting schedule. Truncating a long path frees its slot sooner. Discarding the unfinished answer changes the votes available to admit later requests. To separate these effects, we intervened on the 200 archived GSM8K panels after the initial analysis. The original generation limit was 31,744 tokens. Among 3,200 paths, 550 from 172 questions exceeded 384 tokens, while 146 from 87 questions reached at least 1,024. Clipping their work to 384 while keeping their original votes reduced the mean indexed-over-buffered saving from 8.71\% to 2.62\% at the 3,968-token clock. Keeping the work and replacing only the affected votes by abstentions gave 9.64\%. Applying both changes gave 2.81\%. Thus clipping the long paths accounted for most of the reduction in this archive. The separate interventions use counterfactual terminal votes to isolate each change.

We next compared limits of 384 and 1,024 with the integer parser used on the device. On the same 200 archived questions the modeled mean saving rose from 2.84\% to 7.41\%, a difference of 4.57 percentage points. With the archive parser, the corresponding means were 2.81\% and 7.62\%. We fixed the protocol for a directional test on the A800. We excluded all 200 archived questions and all 16 pilot questions, then selected 32 by a frozen text hash among the 1,105 remaining GSM8K questions. The protocol specifies sixteen reference generations per question, capped at 1,024 tokens, from Qwen2.5-1.5B-Instruct. The short reference truncates these same paths at 384. A path reaching the applicable cap abstained.

At each cap the two gates run with the same fixed path lengths and votes. Four condition orders balance the timed comparisons across questions. For question \(q\) and cap \(c\), let \(s_{q,c}=1-T_{\mathrm{indexed},q,c}/T_{\mathrm{buffered},q,c}\), where \(T\) measures synchronized wall time. Our primary outcome was the question mean of \(s_{q,1024}-s_{q,384}\). The frozen analysis uses a two-sided 95\% bootstrap interval with 20,000 draws of 32 paired questions. The separate practical criterion at 1,024 tokens required at least 2\% saving against buffered admission at that cap, a positive lower confidence limit and at least one changed dispatch.

The short-cap saving was 0.86\% (95\% interval $-0.33$ to 2.23\%) and the long-cap saving was 1.30\% (0.33 to 2.39\%). Their paired difference was $+0.44$ percentage points (95\% interval $-0.79$ to $+1.84$). The trial leaves the directional hypothesis unresolved. The long condition had a positive lower endpoint but fell short of the prespecified 2\% useful saving. Figure~\ref{fig:cap} retains all question-level slowdowns and savings.

Only 35 of the 512 newly sampled paths reached 384 tokens, and two reached 1,024. On 20 of 32 questions the two limits gave identical fixed work and votes. Dispatch changed between gates on five questions at the short cap and seven at the long cap. The indexed gate needed fewer engine steps on three questions at each cap. Replaying the measured paths classified all 64 question--cap dispatch comparisons correctly. Its estimated change in saving was $-0.32$ percentage points, opposite in sign to the wall-time point estimate. Thus correct predictions about which questions take different scheduling actions can coexist with uncertainty about the wall-time difference. The single device session does not separate session variation from a persistent timing effect. Appendix~\ref{app:cap} reports the absolute costs at both output limits.
